\chapter{Division I: short and long division}
\label{ch:div}

Division is harder than multiplication. Multiplying needs no guessing; dividing does. In
long division on paper you look at the leading digits, \emph{guess} the next quotient digit,
multiply back, and correct the guess if it was wrong. Every division algorithm in BigMath is
a refinement of that loop: make the guess cheaper, make it more accurate, or guess many
digits at once.

Given $a$ and $b > 0$, division produces a quotient $q$ and remainder $r$ with
\[
a = q \cdot b + r, \qquad 0 \le r < b .
\]
The tests check exactly this identity for every algorithm and every result.

\section{Dividing by a single limb}
\label{sec:shortdiv}

When the divisor fits in one limb, division is ``short division'': walk from the top limb
down, keeping a running remainder.

\begin{example}[: $9\,876 \div 7$ by short division]
\[
\begin{array}{l@{\qquad}l}
9 \div 7 = 1 \text{ rem } 2 & \\
(2\cdot 10 + 8) = 28 \div 7 = 4 \text{ rem } 0 & \\
(0 \cdot 10 + 7) = 7 \div 7 = 1 \text{ rem } 0 & \\
(0 \cdot 10 + 6) = 6 \div 7 = 0 \text{ rem } 6 & \Rightarrow\ q = 1410,\ r = 6 .
\end{array}
\]
In the library each step divides a two-limb value (remainder, next limb) by the one-limb
divisor.
\end{example}

The step ``divide a two-limb number by a one-limb number'' is the bottleneck: hardware
division of a 128-bit number is slow, or (on ARM) does not exist and becomes a library call.
But the divisor is the same at every step, so BigMath pays for one division up front and
then only multiplies.

\subsection{Dividing by multiplying with a precomputed reciprocal}

\begin{another}[: dividing by 7 without dividing]
To divide many numbers by 7, work out $1/7 \approx 0.14$ once. Then:
\begin{itemize}
  \item $93 \times 0.14 = 13.02$, guess $q = 13$; check $93 - 13 \cdot 7 = 2$, which is
        less than 7, so $q = 13$, $r = 2$.
  \item $98 \times 0.14 = 13.72$, guess $q = 13$; check $98 - 91 = 7$, which is \emph{not} less
        than 7, so correct once: $q = 14$, $r = 0$.
\end{itemize}
The guess may be slightly low, but a cheap check-and-correct step fixes it. Multiplication
replaced division.
\end{another}

The library uses the Möller--Granlund method (2011), which makes this exact for limbs: from
the divisor $d$ it precomputes $v = \floor{(B^2 - 1)/d} - B$. Each step is then one
$64 \times 64 \to 128$ multiplication, a few additions, and at most two small corrections.
Before the precomputation the divisor is \emph{normalized}: shifted left so its top bit is set.
\code{ClassicDivision::DivRem1} keeps the running remainder in that shifted form throughout and
shifts each dividend limb on the fly, so nothing is shifted back and forth on the critical
path.

\begin{perfnote}
Replacing the 128-bit hardware/library division with the reciprocal made printing small
numbers 2.1 times faster (printing divides by $10^{19}$ repeatedly). Keeping the remainder
normalized brought the per-limb cost from 6.4 to 4.9\,ns.
\end{perfnote}

\section{Normalization}
\label{sec:normalize}

Long division needs a good guess for each quotient digit, and the guess is made from the
leading digits. If the divisor's leading digit is tiny, the guess can be far off: dividing by
$10\ldots$ looks the same as dividing by $19\ldots$ when you only see one digit. The fix is to
scale both numbers so the divisor's leading digit is large.

\begin{keyidea}
Shift the dividend and divisor left by the same number of bits so the divisor's top limb has
its highest bit set. The quotient does not change (both were scaled by the same factor), the
remainder comes out scaled and is shifted back at the end. After normalization, a
quotient-digit guess from the leading limbs is never too small and at most two too large
(Knuth's theorem).
\end{keyidea}

In base 10 the equivalent is ``multiply both numbers so the divisor's first digit is at
least 5''.

\section{Long division: Knuth's Algorithm D}
\label{sec:knuth}

For divisors of 2 to 95 limbs BigMath uses classic long division in the form Donald Knuth
wrote down as ``Algorithm D''. Each round:

\begin{enumerate}
  \item \textbf{Guess} the next quotient digit $\hat q$ from the top of the current remainder.
  \item \textbf{Multiply and subtract}: remainder $\mathrel{-}= \hat q \cdot b$, shifted to the
        current position. This is one call of the \code{SubMul1} kernel.
  \item \textbf{Correct}: if the result went negative, the guess was one too big; add $b$
        back and decrement $\hat q$.
\end{enumerate}

\begin{example}[: $7369 \div 47$, the way Algorithm D does it]
\textbf{Normalize.} $47$ starts with 4; multiply both by 2: divide $14738$ by $94$.

\textbf{Round 1.} Window $147$. Guess from the leading digits: $14 \div 9 = 1$.
$147 - 1 \times 94 = 53$. Digit: 1.

\textbf{Round 2.} Bring down 3: window $533$. Guess $53 \div 9 = 5$.
$533 - 5 \times 94 = 63$. Digit: 5.

\textbf{Round 3.} Bring down 8: window $638$. Guess $63 \div 9 = 7$. But
$7 \times 94 = 658 > 638$: the guess was too big. Correct: digit 6, remainder
$638 - 564 = 74$.

\textbf{Un-normalize.} Quotient $156$; remainder $74 \div 2 = 37$.

Check: $47 \times 156 + 37 = 7332 + 37 = 7369$.
\end{example}

\subsection{A better guess: the 3-by-2 quotient digit}

In round 3 above the guess was wrong because it only looked at two digits of the remainder
and one of the divisor. Looking at \emph{three} digits of the remainder and \emph{two} of the
divisor gives $638 \div 94 = 6$: correct immediately. BigMath does exactly that with limbs. It
divides the top three limbs of the remainder by the top two limbs of the divisor, using a
Möller--Granlund \emph{3-by-2} reciprocal precomputed once per division. The guess is then
almost always exact; it is at most one too big, and that happens with probability about
$2/B \approx 10^{-19}$, so the add-back branch practically never runs.

\begin{perfnote}
The library's division loop used to be
\begin{lstlisting}
ULong128 product = (ULong128)v[i] * qhat + borrow;
if (u[j + i] < (DataT)product) { ...borrow... }   // a branch per limb
\end{lstlisting}
with a limb-sized borrow chain and an unpredictable branch on every limb: 2.9--4.0 cycles per
limb. The \code{SubMul1} kernel uses the two-chain layout of \cref{sec:flagchain}: products
along one chain, subtraction along another, no branches: 1.6--1.8 cycles per limb, level with
GMP's assembly from 128 limbs. The routine also works on raw buffers: both operands are
shifted into one workspace (on the stack up to 192 limbs), so a division allocates only its
two results. Divisions with 2- to 256-limb divisors got 1.7 to 2.5 times faster.
\end{perfnote}

\subsection{Cost}

Each quotient limb costs one pass over the divisor, so dividing an $n_a$-limb number by an
$n_b$-limb number costs about $(n_a - n_b + 1) \times n_b$ limb operations. That is fine for a
short divisor or a short quotient, but for two 10\,000-limb numbers with a 10\,000-limb
quotient it is 100 million operations. \Cref{ch:fastdiv} shows how to turn division into a
handful of multiplications instead.

\begin{practice}
For a one-limb divisor, use the scalar overloads (\code{a / uint64\_t}, \code{a \% uint64\_t}):
they go straight to \code{DivRem1}. If you need many divisions by the same multi-limb divisor,
see \code{CachedDivision} in \cref{sec:cached}.
\end{practice}
